\documentclass[book.tex]{subfiles}

\begin{document}
\chapter{Correlations and Renormalization Group Theory}

\label{chap:renorm_group}
\noindent\rule{11cm}{0.4pt} \\
\textbf{Prerequisites:} \autoref{chap:ising_model}, \autoref{chap:ising_critical}  \\
\textbf{Difficulty Level:} *** \\
\noindent\rule{11cm}{0.4pt}
\vspace{5mm}

%\textcolor{red}{TODO: Is this chapter too advanced for now? }

In this chapter we discuss the effect of range of correlations on phase transitions. This correlations can be experimentally measured by scattering experiments, as we will review in subsequent chapters. We also introduce the renormalization group theory. This method has the particularity that it accounts for large length scale fluctuations and is essential for the study of phase transitions. We start with the one-dimensional Ising model, as several of the concepts of the Renormalization Group theory can be illustrated for this case. Then we end with the study of two-dimensional systems.

\section{Correlations in a system}

We showed earlier that the critical point is associated with a divergence in the fluctuations in the magnetization, energy, or number of particles. We call these variables, \textit{order parameters}, \textit{i.e} fluctuating variables whose average values provide the information about the order or broken symmetry of the system. From Monte Carlo simulations, we observe that these fluctuations are highly correlated as we approach the critical point.

It is useful to introduce the concept of \textit{range of correlations}, which consider the distance over which fluctuations in one region  $(s_i - \langle s_i \rangle)$ of space, are correlated or affected by those in another region $(s_j - \langle s_j \rangle)$. Mathematically:

\begin{align}\label{Eq1}
c_{ij} &= \langle (s_i - \langle s_i \rangle)(s_j - \langle s_j \rangle)\rangle \\
&=\langle s_i s_j \rangle -2 \langle s_i \rangle \langle s_j \rangle + \langle s_i \rangle \langle s_j \rangle \\
&= \langle s_i s_j \rangle - \langle s_i\rangle\langle s_j\rangle 
\end{align}

\noindent where $c_{ij}$ is called the \textit{correlation function}. 

If two sites are separated over a larger distance than the range predicted by $c_{ij}$, these two points are uncorrelated of each other. In terms of the magnetic model, two correlated spins are likely to be aligned in the same direction. Thus, we can anticipate that the spontaneous magnetization phenomena is associated with \textit{long range correlations}.

In fact, there is a direct relationship between the correlation function and the fluctuations in magnetization, as we explain bellow:

\begin{align}
\left(\frac{\partial \langle M\rangle}{\partial \beta h}\right)_{\beta,N} &= \langle (\delta M)^2\rangle \\
&= \langle (M - \langle M \rangle)(M - \langle M \rangle)\rangle \\
&= \Bigg\langle \left(\sum_{i=1}^N \left[ s_i - \langle s_i\rangle \right] \right)\left(\sum_{j=1}^N \left[ s_j - \langle s_j\rangle \right]\right)\Bigg\rangle \\
&=\sum_{i=1}^N\sum_{j=1}^N \left[\langle s_i s_j \rangle - \langle s_i\rangle\langle s_j\rangle\right] \\
&=N\sum_{j=1}^N c_{1,j} 
\end{align}

\noindent where in the last part, we consider the fact that the lattice sites are equivalent.

Thus, the divergence of fluctuations at the critical point is directly related with long range of correlations between neighboring spins. When this occurs, the distance of these correlations becomes macroscopic, so even very distant spins will tend to align, forming the big blobs observed in the Monte Carlo simulation.

In the mean-field approximation, we assumed an average magnetization for the neighboring spins, $\langle s_i \rangle = m$ ($c_{i j} = 0$ for $i \neq j$), thus neglecting spatial correlations. This explains the discrepancies of our approximation  with respect to the exact solution. Quantifying the effective interaction between distal spins is a powerful way to think about the physical effect of correlations.

\section{Renormalization Group Theory}

In order to conceptualize the correlations in the Ising model, it is extremely useful to understand how the correlations lead to effective long-range interactions. \textit{Renormalization Group Theory} %\cite{Chandler1987} \cite{Maris1978} \cite{Wilson1971} \cite{Wilson1971a}  
is an analytical method of coarse graining an interacting system to address how physical interactions are impacted by spatial correlations correlations.

The basic procedure of Renormalization Group Theory is:

\begin{itemize}
\item Begin with the original theory with all of relevant degrees of freedom (e.g. the spins $s_i$ in the Ising model).
\item Separate degrees of freedom into one set that will be summed over first and one that will summed over second.
\item Perform the summation over the set of degrees of freedom.
\item Define a new theory based on the remaining degrees of freedom that has coarse-grained (or renormalized) parameters.
\item Identify the renormalization group equations that pass the theory to the next level, successively solve these equations
\end{itemize}

\section{Renormalization Group Theory of the 1D Ising model}
Consider the 1D Ising model with $h = 0$, defined according to the system energy: 
\begin{align}
\beta E=-J\sum_{i=1}^N s_i s_{i+1}
\end{align}
The partition function is defined by:
\begin{align}
Q(J,N) = \sum_{s_1,s_2,...,s_N=-1,1}\exp[J(...+s_1s_2+s_2s_3+s_3s_4++s_4s_5+...)]
\end{align}
Although this has an exact solution, we will employ Renormalization
Group Theory to address the physics of this model (Fig.~\ref{fig1}).

\begin{marginfigure}
\includegraphics[width=\linewidth]{figures/Physics/statistical_mechanics/correlations_and_renormalization_group_theory/renormalized_ising_1d.png}
\caption{Schematic of the renormalization of the 1D Ising model. The 1D Ising model begins with an interaction parameter $J$. The decimation process, where half the spins are summed over, results in an effective 1D Ising model with a new interaction parameter $J'$.}
\label{fig1}
\end{marginfigure}

Arrange the partition into terms that isolate the even spins, \textit{i.e.}:
\begin{align}
Q(J,N) = \sum_{s_1,s_2,...,s_N=-1,1}\exp[J(s_1s_2+s_2s_3)]\exp[J(s_3s_4+s_4s_5)]...
\end{align}
Perform the sum over all even spins, arriving at the expression:
\begin{align}
Q(J,N) = \sum_{s_1,s_3,...,s_N=-1,1}\left[e^{J(s_1+s_3)} + e^{-J(s_1+s_3)}\right]\cdot \\ 
\left[e^{J(s_3+s_5)}+e^{-J(s_3+s_5)}\right]...
\end{align}

Now, recast the partial-summed theory to look the same as the original theory with $N/2$ spins and (perhaps) a new coupling constant $J'$. This is accomplished if we can find a function $f (J)$ and coupling constant $J'$ that satisfies 
\begin{align}
e^{J(s+s')}+e^{-J(s+s')} = f(J)e^{J'ss'}
\end{align} 
for any values of $s, s' = -1,1$. If this transformation is possible, the partition function becomes:

\begin{align}\label{Eq2}
Q(J,N) &= \sum_{s_1,s_2,...,s_N=-1,1} f(J) e^{J's_1s_3}f(J) e^{J's_3s_5}... \\
&= [f(J)]^{N/2}Q(J',N/2)]
\end{align}

\noindent where $Q(J',N/2)$ has the same form as the original theory.

To find $f(J)$ and $J'$, consider two scenarios:
\begin{itemize}
\item If $s = s' = \pm 1$, the coarse-graining criterion is:

\begin{align}\label{Eq3}
e^{2J}+e^{-2J}=f(J)e^{J'}
\end{align}
\item If $s = -s' = \pm 1$, the coarse-graining criterion is:
\begin{align}\label{Eq4}
2=f(J)e^{-J'}
\end{align}
\end{itemize}

These two criteria are easily solved to get:
\begin{align}
J' &= \frac{1}{2} \log [\cosh(2J)]\\
f(J) &= 2\cosh^{1/2}(2J) 
\end{align}

\noindent and we define $g(J) = N^{-1} \log Q$, resulting in the expression:
\begin{align}
g(J')=2g(J)-\log\left[ 2\sqrt{\cosh(2J)} \right] 
\end{align}

\noindent These are the \textit{renormalization group equations}. The inverses of these equations are:

\begin{align}
J&=\frac{1}{2}\cosh^{-1}\left( e^{2J'} \right) \hspace{0.2cm} \\
g(J) &= \frac{1}{2} g(J') + \frac{1}{2} \log 2 + \frac{J'}{2} 
\end{align}

Successive application of the RG equations always decreases the coupling constant $J$, except for the initial value $J \rightarrow \infty$ (Fig.~\ref{fig2}). The loss of coupling implies the lack of long-range order associated with a phase transition. Conditions where $J' = J$ is called a fixed point of the theory, which identifies a state whose physical behavior is length-scale invariant (\textit{self-similar}). The 1D Ising model has a stable fixed point at $J = 0$ and an unstable fixed point at $J \rightarrow \infty$.

\begin{marginfigure}
\includegraphics[width=\linewidth]{figures/Physics/statistical_mechanics/correlations_and_renormalization_group_theory/rg_equation_coupling.png}
\caption{Value of the coupling parameter with number of applications
of the RG equations}
\label{fig2}
\end{marginfigure}

\section{Renormalization Group Theory of the 2D Ising model}

The 2D Ising model can also be addressed using Renormalization Group theory. We follow a similar procedure as in the 1D case, though several approximations need to be made.

The partition function organized into terms such that:

\begin{align}
Q = \sum_{s_1,s_2,...,s_N=-1,1} \dotsc \exp[Js_5(s_1+s_2+s_3+s_4)]  \\ 
\times \exp[Js_6(s_2+s_3+s_7+s_8)] \dotsc
\end{align}

\noindent where $s_5$ and $s_6$ are among the spins that we will first sum over (Fig~\ref{fig3}).

\begin{marginfigure}
\includegraphics[width=\linewidth]{figures/Physics/statistical_mechanics/correlations_and_renormalization_group_theory/renormalization_2d_coupling.png}
\caption{Schematic of the renormalization of the 2D Ising model}
\label{fig3}
\end{marginfigure}

Sum over half the spins to get:
\begin{align}
Q = \sum_{\text{remaining }s_i} \dotsc \left[ e^{J(s_1+s_2+s_3+s_4)} + e^{-J(s_1+s_2+s_3+s_4)} \right] \\
\times \left[ e^{J(s_2+s_3+s_7+s_8)} + e^{-J(s_2+s_3+s_7+s_8)} \right] \dotsc
\end{align}

It would be desirable to find a set of renormalization group equations based on:
\begin{align}\label{Eq5}
e^{J(s_1+s_2+s_3+s_4)} + e^{-J(s_1+s_2+s_3+s_4)} = f(J) e^{J'(s_1s_2+s_1s_4+s_2s_3+s_3s_4)}
\end{align}

\noindent however this will prove to be impossible.

Since there exist unique scenarios for $s_1$, $s_2$, $s_3$, $s_4$, such that:
\begin{align}
s_1 &=s_2=s_3=s_4 = \pm 1 \\
s_1 &=s_2=s_3=-s_4 = \pm 1 \\
s_1 &=s_2=-s_3=-s_4 = \pm 1 \\
s_1 &=-s_2=-s_3=-s_4 = \pm 1 
\end{align}

\noindent the model is overspecified (4 equations, 2 unknowns), and the coarse-grained theory cannot have the same form as the original.

The simplest possibility is:
\begin{align}
&e^{J(s_1+s_2+s_3+s_4)} + e^{-J(s_1+s_2+s_3+s_4)} = \\
 &f(J)\exp \left[ \frac{1}{2}J_1(s_1s_2+s_2s_3+s_3s_4+s_4s_1) + J_2(s_1s_3+s_2s_4) + J_3s_1s_2s_3s_4 \right] 
\end{align}

Since there are 4 equations and 4 unknowns, these can be solved to find:

\begin{align}\label{Eq6}
J_1 &= \frac{1}{4}\log[\cosh(4J)] \\
J_2 &= \frac{1}{8}\log[\cosh(4J)] \\
J_3 &= \frac{1}{8}\log[\cosh(4J)] - \frac{1}{2}\log[\cosh(4J)] \\
f(J) &= 2[\cosh(2J)]^{1/2}[\cosh(4J)]^{1/8}
\end{align}

Inserting this into the partition function, we have:

\begin{align}\label{Eq7}
Q(J,N) = [f(J)]^{N/2} \sum_{N/2 spins}\exp\left[ J_1\sum_{ij}{'}s_i s_j + \right. \\
 \left. J_2\sum_{lm}{''}s_l s_m + J_3\sum_{pqrt}{'''}s_p s_q s_r s_t \right] 
\end{align}

\noindent where $\sum{'}$ is over nearest neighbors, $\sum{''}$ is over next-nearest neighbors, and $\sum{'''}$ is over four spins in a square.

This demonstrates that coarse-graining introduces new interactions into the theory that did not exist in the original theory. However, we cannot use this structure, since the next renormalization step will introduce even more complex interactions. Therefore, we make the approximation:

\begin{align}\label{Eq8}
J_1\sum_{ij}{'}s_is_j+J_2\sum_{lm}{''}s_ls_m \approx J'(J_1,J_2)\sum_{ij}{'}s_is_j
\end{align}

\noindent and neglect the 4-spin interaction ($J_3 \approx 0$).

This approximation gives:

\begin{align}\label{Eq9}
Q(J,N) = [f(J)]^{N/2} Q[J'(J_1,J_2),N/2]
\end{align}

\noindent and the free energy is found from $g(J)=\frac{1}{2}\log f(J)+\frac{1}{2}g(J')$.

To estimate $J'$, consider the case where all the spins are aligned. Since we have $N/2$, there are $N$ near neighbors and $N$ next-near neighbors, thus $J_1\sum_{i j}{'}s_is_j = NJ_1$ and $J_2\sum_{lm}{''}s_ls_m = NJ_2$. Therefore, we estimate that $J'\approx J_1+J_2$.

Using our RG equations for $J_1$ and $J_2$, we have:
\begin{align}
J'=\frac{3}{8}\log[\cosh(4J)] 
\end{align}

Unlike the 1D Ising model, the RG equation for the 2D system has a non-trivial fixed point $J_c = 0.50698$ that satisfies $J_c = \frac{3}{8}\log[\cosh(4J_c)]$. The critical point $J_c$ dictates conditions where the correlations diverge at all length scales (i.e. the system is self-similar) and marks the onset of a phase transition (Fig.~\ref{fig4}).

\begin{marginfigure}
\includegraphics[width=\linewidth]{figures/Physics/statistical_mechanics/correlations_and_renormalization_group_theory/coupling_constant_values.png}
\caption{Value of the coupling parameter with number of applications
of the RG equations}
\label{fig4}
\end{marginfigure}

\newpage

\section{Coarse-Grained Physical Models}

Renormalization Group theory gives a systematic method for coarse graining a physical model. In many instances, the insights that can be gained by coarsegraining can provide fundamental insight into the dominant physics. For example, coarse-graining concepts have been pivotal in polymer physics (Fig.~\ref{fig5},~\ref{fig6}).

\begin{marginfigure}
\includegraphics[width=\linewidth]{figures/Physics/statistical_mechanics/correlations_and_renormalization_group_theory/Fig5.png}
\caption{\textbf{Discretization of a polymer chain into equivalent effective representations.} This schematic demonstrates the process of renormalization of the polymer chain by summing over intermediate degrees of freedom. Each subsequent step results in effective elastic couplings between chain segments that capture the underlying physics.}
\label{fig5}
\end{marginfigure}

\begin{marginfigure}
\includegraphics[width=\linewidth]{figures/Physics/statistical_mechanics/correlations_and_renormalization_group_theory/Fig6.png}
\caption{The blob argument for the behavior of a real chain in a tube results in a picture of the polymer chain feeling a different dimensionality depending on the length scale that is considered. At small length scales, the chain segments live in 3 dimensions. However, the confinement cordons the chain into effective units or blobs that do not penetrate each other due to excluded volume. At larger length scales, these blobs effectively live in 1 dimension due to the confinement.}
\label{fig6}
\end{marginfigure}

\end{document}
